% Journal of Forensic Sciences -- Supplemental Information (upload as a
% separate file, designated "Supplemental Information"). Anonymised.
% Overleaf: set this file as the main document to compile it.
\documentclass[12pt]{article}
% ---------------------------------------------------------------- packages
\usepackage[T1]{fontenc}
\usepackage[utf8]{inputenc}
\usepackage{mathptmx}                 % Times text and maths
\usepackage[scaled=0.92]{helvet}      % Helvetica (line numbers)
\usepackage{courier}                  % Courier (file names)
\usepackage[letterpaper,margin=1in]{geometry}
\usepackage{amsmath,amssymb}
\usepackage{graphicx}
\usepackage[table]{xcolor}
\usepackage{booktabs,array,makecell}
\usepackage{adjustbox}
\usepackage{pifont}
\usepackage{algorithm}
\usepackage{algpseudocode}
\usepackage{tcolorbox}
\usepackage{float}
\usepackage{etoolbox}
\usepackage{setspace}
\usepackage{caption}
\usepackage{titlesec}
\usepackage[numbers,square,sort&compress]{natbib}
\usepackage{lineno}
\usepackage{xurl}
\urlstyle{same}
\usepackage[hidelinks]{hyperref}
\hypersetup{pdftitle={Structure-aware forgery detection for Aadhaar and Permanent Account Number cards},
  pdfauthor={}, pdfsubject={}, pdfkeywords={}, pdfcreator={}}

% ---------------------------------------------------------------- JFS layout
% Double spacing throughout, including abstract, highlights, references
% and table and figure legends (JFS Information for Authors, June 2025).
\doublespacing
\setlength{\parindent}{1.5em}
\pagestyle{plain}

% Headings: all caps and bold; subheadings upper/lower case, italic, not
% bold; sub-subheadings upper/lower case, plain text.
\setcounter{secnumdepth}{3}
\titleformat{\section}{\normalfont\bfseries}{\thesection}{0.6em}{\MakeUppercase}
\titleformat{\subsection}{\normalfont\itshape}{\thesubsection}{0.6em}{}
\titleformat{\subsubsection}{\normalfont}{\thesubsubsection}{0.6em}{}
\titleformat{\paragraph}[runin]{\normalfont\itshape}{}{0pt}{}
\titlespacing*{\section}{0pt}{1.2\baselineskip}{0.4\baselineskip}
\titlespacing*{\subsection}{0pt}{0.9\baselineskip}{0.3\baselineskip}
\titlespacing*{\subsubsection}{0pt}{0.7\baselineskip}{0.2\baselineskip}
\titlespacing*{\paragraph}{0pt}{0.4\baselineskip}{0.5em}

% Legends: "TABLE 1 Legend text." and "FIGURE 1 Legend text."
\captionsetup{font=doublespacing, labelfont=bf, labelsep=space,
  justification=justified, singlelinecheck=false, skip=6pt}
\captionsetup[table]{name=TABLE, position=top}
\captionsetup[figure]{name=FIGURE}

% Tables: single-spaced body, 10 pt; each tabular is shrunk only if it is
% wider than the text (adjustbox, max width).
\AtBeginEnvironment{table}{\let\scriptsize\footnotesize\let\tiny\scriptsize}

% References: bold "References" heading, numbered "1.", "2.", ...
\renewcommand{\bibsection}{\par\bigskip\noindent\textbf{References}\par}
\renewcommand{\bibnumfmt}[1]{#1.}
\setlength{\bibsep}{0pt}

% ---------------------------------------------------------------- paper macros
\algtext*{EndFor}% omit "end for" / "end if" lines in the algorithms
\algtext*{EndIf}

\definecolor{cvis}{HTML}{2F6DB5}    % visual branch
\definecolor{crule}{HTML}{D27A1E}   % structural / semantic branch
\definecolor{cfuse}{HTML}{2E8B57}   % fusion
\definecolor{cbad}{HTML}{C0392B}    % failure / forgery
\definecolor{cink}{HTML}{2B2B2B}
\definecolor{cgrey}{HTML}{6B7280}

\newcommand{\cmark}{\textcolor{cfuse}{\ding{51}}}
\newcommand{\xmark}{\textcolor{cbad}{\ding{55}}}

% Research-question box (Section 1)
\newtcolorbox{rqbox}[1]{colback=white, colframe=black!60,
  boxrule=0.5pt, arc=0pt, left=8pt, right=8pt, top=5pt, bottom=5pt,
  fonttitle=\bfseries\normalsize, coltitle=black, colbacktitle=black!8,
  title=#1, fontupper=\normalsize}

% Auto-generated by scripts/export_latex.py -- do not edit by hand.
\newcommand{\dsPersons}{300}
\newcommand{\dsDocuments}{3{,}600}
\newcommand{\dsImages}{10{,}800}
\newcommand{\dsGenuine}{3{,}600}
\newcommand{\dsForged}{7{,}200}
\newcommand{\dsGiB}{0.88}
\newcommand{\dsNq}{3}
\newcommand{\dsTrainImg}{7{,}560}
\newcommand{\dsValImg}{1{,}620}
\newcommand{\dsTestImg}{1{,}620}
\newcommand{\dsTrainPer}{210}
\newcommand{\dsValPer}{45}
\newcommand{\dsTestPer}{45}
\newcommand{\dsCellImages}{900}
\newcommand{\dsCellControls}{1{,}800}
\newcommand{\dsPhashBits}{64}
\newcommand{\dsPhashDistinctSixtyFour}{5{,}848}
\newcommand{\dsPhashDistinct}{5{,}848}
\newcommand{\dsPhashAUC}{0.914}
\newcommand{\dsPhashOverlap}{9.9}
\newcommand{\dsPhashWithinMed}{6}
\newcommand{\dsPhashCrossPersonMed}{24}
\newcommand{\dsPhashCrossPersonMin}{0}
\newcommand{\dsAadhaarSpace}{8\times10^{10}}
\newcommand{\dsPanSpace}{4.6\times10^{10}}
\newcommand{\dsAadhaarStrings}{1{,}083}
\newcommand{\dsPanStrings}{1{,}200}
\newcommand{\dsAadhaarCoinPct}{1.8}
\newcommand{\dsAadhaarCoinExp}{19}
\newcommand{\dsPanCoinPct}{1.8}
\newcommand{\dsPanCoinExp}{21}
\newcommand{\dsStrictAll}{9.0}
\newcommand{\dsStrictInitialFirst}{29.3}
\newcommand{\dsStrictLeadingInitial}{93.1}
\newcommand{\dsCICell}{2.0}
\newcommand{\dsCIPerson}{3.4}
\newcommand{\dsFPAadhaarClean}{0.0}
\newcommand{\dsFPAadhaarMild}{0.0}
\newcommand{\dsFPAadhaarSevere}{11.5}
\newcommand{\dsFPAadhaarAll}{3.8}
\newcommand{\dsFPPanClean}{13.7}
\newcommand{\dsFPPanMild}{12.0}
\newcommand{\dsFPPanSevere}{33.3}
\newcommand{\dsFPPanAll}{19.7}
\newcommand{\dsOcrAadhaarAadhaarNumberClean}{99.9}
\newcommand{\dsOcrAadhaarAadhaarNumberMild}{99.8}
\newcommand{\dsOcrAadhaarAadhaarNumberSevere}{88.3}
\newcommand{\dsOcrPanPanNumberClean}{84.5}
\newcommand{\dsOcrPanPanNumberMild}{85.3}
\newcommand{\dsOcrPanPanNumberSevere}{66.7}
\newcommand{\dsOcrAadhaarNameClean}{98.9}
\newcommand{\dsOcrAadhaarNameMild}{98.7}
\newcommand{\dsOcrAadhaarNameSevere}{64.4}
\newcommand{\dsOcrPanNameClean}{99.8}
\newcommand{\dsOcrPanNameMild}{99.3}
\newcommand{\dsOcrPanNameSevere}{54.6}
\newcommand{\dsProbeRecords}{40}
\newcommand{\dsProbeCells}{8}
\newcommand{\dsProbeComparisons}{320}
\newcommand{\dsProbeOutside}{0}
\newcommand{\dsProbeCritical}{0.500}
\newcommand{\dsProbeCriticalAUC}{0.493}
\newcommand{\dsProbeCriticalN}{1{,}080}
\newcommand{\dsProbeCriticalCI}{$[0.470,\,0.530]$}
\newcommand{\dsProbeControl}{0.491}
\newcommand{\dsProbeControlAUC}{0.494}
\newcommand{\dsProbeControlN}{540}
\newcommand{\dsProbeControlCI}{$[0.449,\,0.533]$}
\newcommand{\dsProbeCriticalClean}{0.492}
\newcommand{\dsProbeCriticalCleanAUC}{0.491}
\newcommand{\dsProbeCriticalCleanN}{360}
\newcommand{\dsProbeCriticalCleanCI}{$[0.440,\,0.543]$}
\newcommand{\dsProbeControlClean}{0.506}
\newcommand{\dsProbeControlCleanAUC}{0.454}
\newcommand{\dsProbeControlCleanN}{180}
\newcommand{\dsProbeControlCleanCI}{$[0.433,\,0.579]$}


\renewcommand{\thetable}{S\arabic{table}}
\onehalfspacing
\begin{document}
\begin{center}
{\large\bfseries Supplemental Information\par}
\medskip
Structure-Aware Forgery Detection for Aadhaar and Permanent Account Number Cards: A Controlled Study of Visual, Structural, and Semantic Evidence Fusion
\end{center}
\bigskip
\section*{Appendix S1. Reproducibility Record}



This appendix states what a reader needs to repeat each stage of the
study, how exactly each stage can be repeated, and every change made to
the protocol after it was approved. The repository file
\texttt{docs/REPRODUCIBILITY.md} gives the commands and the SHA-256
digest of every frozen file.

\subsection*{Levels of Reproduction}


Table~\ref{tab:repro-levels} separates four levels, because they do not
reproduce equally well.

\begin{table}[H]
\caption{What can be repeated, and how exactly.}
\label{tab:repro-levels}
\centering
\scriptsize
\setlength{\tabcolsep}{3pt}
\renewcommand{\arraystretch}{1.12}
\begin{adjustbox}{max width=\linewidth}
\begin{tabular}{@{}p{0.26\columnwidth} p{0.31\columnwidth} p{0.35\columnwidth}@{}}
\toprule
Level & Starting point & Agreement observed \\
\midrule
L1 Scoring & released per-image scores, frozen calibration and fusion files
  & exact; checked on validation by CI at every commit \\
L2 Model inference & released checkpoints and corpus
  & logits equal to the 6-decimal precision of the score files
    (max.\ difference $5{\times}10^{-7}$, 40 images per seed) \\
L3 Corpus & generator, configuration and seed
  & 10,791 of 10,800 images byte-identical on another OS; 9 differ in ${\le}37$ pixels by ${\le}3$ levels \\
L4 Training & training script, corpus, seed
  & bit-identical on the same GPU and library versions (one seed re-run); other hardware not tested \\
\bottomrule
\end{tabular}
\end{adjustbox}
\end{table}

\emph{L1.} Every number in Section~4 of the main text is computed by
committed scripts from per-image prediction files, and its tables are
generated from their output. A continuous-integration job
re-applies the frozen calibration and fusion files to the validation
split and fails if the rounded summary changes by one character. The
test and fresh-set summaries are produced by the same code.

\emph{L2.} Before the fresh set was scored, each of the five tile
checkpoints re-scored 40 v1 test images and matched its stored logits
to the stored precision. Inference is therefore deterministic for
practical purposes on the same hardware.

\emph{L3.} The corpus was built under Windows. We rebuilt all 10,800
images under Linux with the same pinned library versions. The first
attempt differed in almost every image: the imaging library applies
complex-script text layout only when an optional library is installed,
which is usually the case on Linux and not on Windows, so text landed a
few pixels apart. With the generator now fixed to the plain layout
engine on every platform, 10,791 images are byte-identical. The other
nine differ in at most 37 pixels, by at most 3 of 255 grey levels, all
at the right image border, which is consistent with a floating-point
difference in the perspective warp. Labels, identifiers, splits and
ground-truth rule verdicts are identical. Per-image digests are
released so that any rebuild can be checked image by image. One
consequence of the plain layout engine is that Devanagari text is drawn
without complex-script shaping; no Devanagari field is validated.

\emph{L4.} The seed fixes the weight initialisation of the new head,
the data order and the negative tiles drawn each epoch, and cuDNN
auto-tuning is off. Deterministic GPU kernels were not enforced, and
training uses half precision, so before checking we expected a second
run of the same seed to differ slightly. To check, tile seed~0 was
trained again from a fresh copy of the repository, in a new Python
environment, on the same GPU with the same Python, PyTorch and CUDA
versions (the OS build had changed from 26200 to 26300). The pass rule
was recorded before the run: identical training configuration,
identical input files, and best validation AUROC and PR-AUC inside the
range of the five original seeds (.685--.773 and .611--.684). Under
this rule the run formally \emph{failed}. The recorded digest of the
edit-region file differed because Git for Windows had converted the
file's line endings on checkout; with the conversion undone the file is
identical to the original, and the repository is now set so that Git
converts no file. The metric check passed. The re-run also matched the
original more closely than the rule required: the loss and validation
metrics of every epoch, the per-image training and validation scores,
and the selected checkpoint are bit-identical to the original run
(validation AUROC .699, PR-AUC .627, best epoch 12 of 17 in both). On
the same hardware and library versions, training is therefore exactly
repeatable. Other GPUs, drivers or library versions were not tested,
and half-precision kernels can differ between them, so we still treat
the five seeds, not a single run, as the unit of variation
(Table~\ref{tab:runs}).

\subsection*{Training Runs and Environment}

Table~\ref{tab:runs} lists every visual model trained, including the
two variants that were not selected. All runs used one NVIDIA GeForce
RTX 4060 Laptop GPU (8\,GB), Python 3.12.7, PyTorch 2.14.1,
torchvision 0.29.1, CUDA 13.0 and cuDNN 9.24 (driver 616.92, read
after the runs), under Windows~11 (build
26200; whole-card seeds 3--4 and the fresh-set scoring ran on build
26300 after an automatic update, as did the re-run below). Training
took 7.3 GPU-hours in total (3.7 for the tile model), not counting the
re-run. Each run records the code commit it ran from
and refuses to start if tracked files are modified. The runs span six
commits; we checked that the tile training script is identical across
its runs, and that the whole-card script changed only by adding the
variant switch used for the regularised pilot.

\begin{table}[H]
\caption{All visual training runs (validation, C0 vs.\ C1). Ep.: best
epoch / epochs run; $T$: fitted temperature; ECE after scaling.
Re-run: clean re-run of seed~0 (L4, Appendix~S1); used in no
result and not calibrated.}
\label{tab:runs}
\centering
\scriptsize
\setlength{\tabcolsep}{2.6pt}
\renewcommand{\arraystretch}{1.08}
\begin{adjustbox}{max width=\linewidth}
\begin{tabular}{@{}l c c r c c c c@{}}
\toprule
Model & Seed & Ep. & Min & AUROC & PR-AUC & $T$ & ECE \\
\midrule
Tiles & 0 & 12/17 & 42.6 & .699 & .627 & 4.40 & .109 \\
      & 1 & \hphantom{0}6/11 & 30.9 & .685 & .611 & 1.52 & .055 \\
      & 2 & 17/22 & 53.9 & .735 & .661 & 4.54 & .052 \\
      & 3 & 15/20 & 49.3 & .717 & .661 & 4.16 & .089 \\
      & 4 & 13/18 & 45.0 & .773 & .684 & 3.42 & .058 \\
      & ens. & -- & -- & \textbf{.780} & \textbf{.702} & -- & .073 \\
      & 0, re-run & 12/17 & 42.3 & .699 & .627 & -- & -- \\
\addlinespace[4pt]
Whole card & 0 & 18/23 & 38.6 & .578 & .448 & 6.48 & .098 \\
           & 1 & 20/25 & 41.6 & .613 & .500 & 5.40 & .082 \\
           & 2 & 19/24 & 39.9 & .629 & .501 & 4.94 & .052 \\
           & 3 & 18/23 & 39.9 & .622 & .507 & 5.17 & .078 \\
           & 4 & 18/23 & 37.9 & .611 & .518 & 5.90 & .112 \\
           & ens. & -- & -- & .638 & .531 & -- & .082 \\
\addlinespace[4pt]
Regularised & 0 & \hphantom{0}5/10 & 18.5 & .564 & .435 & -- & -- \\
\bottomrule
\end{tabular}
\end{adjustbox}
\end{table}

\subsection*{Protocol Record}

Table~\ref{tab:protocol} lists, in order, every change to the approved
protocol and every irregular event during the runs. Each entry is dated
in the protocol file; entries made after the test split was scored are
marked as such. Two points need to be stated explicitly. First, the
training script wrote the test-split score files when each model
finished and recorded their digests; the files were not read before the
freeze, and the scoring script refuses to run unless the freeze tag
exists, the frozen files equal their tagged versions and every score
file matches its recorded digest. This shows that the files were not changed, not
that they were never opened; the latter rests on the authors' record.
Second, per-check fusion (A5) was designed after the test split had
been scored, so on that split it is a post-hoc result. Only its
confirmation on the fresh set, which was generated after A5 was fixed,
is confirmatory.

\begin{table}[H]
\caption{Protocol record: amendments, corrections and run events, in
order. ``Test'' gives the state of the test split when the entry was
made.}
\label{tab:protocol}
\centering
\fontsize{8.5}{10}\selectfont
\setlength{\tabcolsep}{3.5pt}
\renewcommand{\arraystretch}{1.12}
\begin{adjustbox}{max width=\linewidth}
\begin{tabular}{@{}l l l p{0.40\textwidth} p{0.27\textwidth}@{}}
\toprule
Entry & Date & Test & Change and reason & Effect on reported results \\
\midrule
A1 & before training & sealed
  & Whole-card input raised from $384{\times}256$ to $768{\times}512$; fine-trace and coarse C1 edits reported separately. Reason: edit-survival measurement on training pairs (Fig.~10 in the main text).
  & Defines the comparison model; no score existed. \\
Batch rule & 17 Oct & sealed
  & Batch chosen so that peak reserved memory stays within 92\% of physical VRAM, not merely ``no out-of-memory error'': the Windows driver silently spilled into system RAM.
  & Whole-card batch 16 with 2-step accumulation (effective 32). \\
A2 & 18 Oct & sealed
  & Whole-card model memorised its training set (val.\ AUROC .578). Validation-only development allowed; a regularised variant and a tile model specified, each with its selection rule, before being run.
  & All candidates reported (Table~\ref{tab:runs}). \\
A3 & 18 Oct & sealed
  & Tile model selected (val.\ PR-AUC .627 against a pre-set bar of .478). Planned $384{\times}256$ training arm dropped; resolution analysis replaced by the tiles vs.\ whole-card comparison.
  & Tile model is the visual branch everywhere. \\
Run event & 19 Oct & sealed
  & Windows update stopped whole-card seed 3 in its first epoch; re-run from scratch. OS build changed to 26300 for whole-card seeds 3--4.
  & None expected; libraries and code unchanged. \\
A4 & 19--20 Oct & sealed
  & Fusion rule, features and thresholds fixed on validation; approved by all authors; freeze tag created (commit identifier withheld for double-blind review).
  & Primary pre-registered fusion. \\
Correction & 20 Oct & sealed
  & Recorded corpus digest belonged to a pre-freeze build of 29 Aug; replaced by the digest of the frozen images after a file-by-file check.
  & None; no image, label or score changed. \\
Test scored & 21 Oct & scored once & All frozen systems scored on the test split. & -- \\
A5 & 21 Oct & \textbf{scored}
  & Per-check fusion designed after inspecting test results; design recorded before fitting. Agreed by one author on 21 Oct and confirmed by the other authors on 23 Oct.
  & Reported as post-hoc on the test split. \\
A6 & 21 Oct & scored
  & Fresh 90-person test set pre-registered with hypothesis and analysis, then generated. First build stopped on a text-encoding error (Windows code page) and was re-run with UTF-8; nothing had been scored.
  & Confirmatory test of A5 against A4. \\
OCR record & 21 Oct & scored
  & Protocol listed Tesseract 5.3.4; the engine that read v1 and the fresh set was 5.5.3 (Leptonica 1.87.0), installed before v1 was built.
  & None; one engine produced all OCR text. Record corrected. \\
Generator fix & 24 Oct & scored
  & Text layout fixed to the plain engine on all platforms, the engine that produced the frozen images (Windows); found by an independent Linux rebuild.
  & None; v1 and fresh images unchanged. Rebuilds elsewhere now match (L3). \\
Analysis plan & 28 Oct & scored
  & Families, paired bootstrap and Holm correction for Section~4 of the main text recorded; the comparisons come from lists written before the test run, the grouping was recorded after the aggregate test reports existed.
  & Defines Table~15 of the main text; no prediction or parameter changed. \\
Re-run & 29 Oct & scored
  & Clean re-run of tile seed 0 under a pass rule recorded before the run (L4, Appendix~S1). Formal fail: the digest of the edit-region file changed because line endings were converted on checkout; its content is identical. Metric check passed; run bit-identical to the original.
  & None. Git now set to convert no file on checkout. \\
\bottomrule
\end{tabular}
\end{adjustbox}
\end{table}
\end{document}
